\chapter{Literature Review}

\section{Introduction}

Financial sentiment analysis sits at the intersection of behavioural finance, Natural Language Processing (NLP), and decision-support systems. Research consistently shows that public sentiment from news articles, social media, and online communities influences asset prices, volatility, and trading behaviour, with strong empirical links to short-term market movements \citep{tetlock2007giving, bollen2011twitter}.

This chapter reviews five areas of the literature that directly shaped my design decisions: (i)~investor sentiment and behavioural finance; (ii)~NLP and sentiment analysis in finance; (iii)~Aspect-Based Sentiment Analysis (ABSA); (iv)~explainable AI (XAI) in financial systems; and (v)~sentiment--price relationships and causal inference. The final section pulls these together to identify the gaps that motivated the specific contributions of this project.

\section{Investor Sentiment and Financial Decision Making}

Behavioural finance challenges the assumption of fully rational markets, demonstrating that investor sentiment and psychological factors can produce temporary mispricings \citep{barberis2003behavioral, delong1990noise}.  The efficient market hypothesis \citep{fama1970efficient} would predict that irrational sentiment is quickly arbitraged away, but the empirical record tells a more complicated story.  Tetlock's landmark study \citep{tetlock2007giving} showed that negative language in mainstream financial news correlates with short-term downward pressure on stock returns, with subsequent reversion as markets correct.  Earlier work by Antweiler and Frank \citep{antweiler2004noise} on Yahoo Finance and Raging Bull message boards had already established that user-generated online text contains statistically significant information about market volatility, even if the return-prediction effect is economically modest.  Together these studies established a precedent for using textual data to proxy for market-moving information flows.

Bollen et al.\ \citep{bollen2011twitter} extended this analysis to social media at scale, constructing daily public mood indices from millions of Twitter posts across dimensions including \textit{Calm}, \textit{Happy}, and \textit{Anxious}, and demonstrating statistically significant predictive relationships with subsequent Dow Jones Industrial Average movements.  Their work suggested that emotional tone in online discourse contains forward-looking market information beyond what is captured by price or volume data alone.

More recent studies have expanded the scope of analysis to retail investor communities.  Long et al.\ \citep{long2023reddit} examined Reddit's r/WallStreetBets in the context of the GameStop rally, demonstrating how community-driven sentiment surges can dramatically amplify retail trading volume and short-term price volatility.  Research in cryptocurrency markets similarly finds that sentiment extracted from Twitter \citep{kraaijeveld2020predictive} and even blockchain transactional data \citep{kleitsikas2025bitcoin} can reflect speculative behaviour and anticipate short-term price swings.  Cryptocurrency markets are particularly sensitive to retail sentiment, making multi-source coverage especially important for any system targeting this asset class.

These findings underscored for me the importance of looking at multiple platforms simultaneously. Retail sentiment on Reddit behaves quite differently from institutional-oriented news, and treating either in isolation gives a distorted picture. That was the core motivation for building a multi-source system rather than focusing on a single channel.

\section{NLP and Sentiment Analysis in Finance}

\subsection*{Lexicon-Based and Classical Machine Learning Approaches}

Early financial sentiment systems relied on domain-specific lexicons.  Loughran and McDonald \citep{loughran2011liability} demonstrated that general-purpose word lists (such as the Harvard-IV-4 dictionary) are poorly suited to financial text, because words like ``liability'' or ``depreciation'' carry no negative connotation in that domain.  Their finance-specific dictionary, which assigns positive and negative polarities to vocabulary representative of financial corpora, became the standard lexicon baseline for the field.  While transparent and computationally inexpensive, lexicon-based methods struggle with financial nuance, negation, and context-dependent phrasing.  Classical machine learning approaches applying Naïve Bayes and Support Vector Machines over bag-of-words features improved upon lexicons on structured datasets \citep{kogan2009predicting} but remained limited in their capacity to model semantic relationships across tokens.

These traditional approaches fail on phrases characteristic of financial text, such as ``missed by a whisker'' or ``headwinds persist'', where correct interpretation depends on domain context and sentence structure.  An intermediate generation of models used dense word embeddings \citep{mikolov2013word2vec, pennington2014glove} to capture semantic similarity, and recurrent architectures such as LSTMs \citep{hochreiter1997lstm} to model sequential context.  While these represented a significant improvement over bag-of-words features, they still struggled with long-range dependencies and the absence of bidirectional context.  The broader challenge of extracting sentiment from unstructured text at scale is surveyed in \citeauthor{cambria2016affective} (\citeyear{cambria2016affective}), who frames sentiment analysis as a core component of affective computing with applications spanning from social media monitoring to financial decision support.

\subsection*{Transformer Models and Domain-Specific Pre-training}

The introduction of the transformer architecture \citep{vaswani2017attention} fundamentally changed the landscape of NLP by replacing recurrence with a self-attention mechanism that models dependencies across the full input sequence.  Applying this to language understanding, Devlin et al.\ \citep{devlin2019bert} introduced BERT, which demonstrated that deep bidirectional pre-training followed by task-specific fine-tuning consistently outperformed task-specific architectures across a wide range of NLP benchmarks.  Although general-purpose models such as BERT and RoBERTa \citep{liu2019roberta} perform strongly on broad NLP benchmarks, they exhibit domain-mismatch degradation when applied to financial text, where vocabulary, discourse structure, and polarity cues differ substantially from general corpora.

FinBERT \citep{araci2019finbert} addressed this by domain-adaptively pre-training on large-scale financial corpora, including financial news and analyst reports.  FinBERT consistently outperforms general-purpose models on the Financial PhraseBank benchmark \citep{malo2014financialphrasebank}, achieving accuracy and F1-scores substantially above lexicon baselines.  Its ability to interpret ambiguous phrasing, negation, and implicit sentiment has made it the de-facto standard model for finance-specific sentiment analysis in the academic literature.

More recent large language models targeting finance, including FinGPT \citep{weng2023fingpt} and FinSentGPT \citep{mahda2024finsentgpt}, offer broader inference capacity and multilingual coverage, but their computational cost and limited interpretability make them less practical for academic prototypes or latency-sensitive production environments compared to FinBERT.

\subsection*{Datasets and Evaluation Practices}

Evaluation of financial sentiment models commonly relies on the Financial PhraseBank \citep{malo2014financialphrasebank}, a dataset of approximately 4,840 sentences annotated by financial professionals.  While widely used, it is limited in both size and domain diversity.  Broader datasets drawn from earnings calls, news archives, and social media \citep{schumaker2009textmining, smailovic2014predictive} offer richer coverage but differ in annotation methodology, class balance, and noise characteristics.  These discrepancies pointed me toward building my own evaluation dataset rather than relying on an existing benchmark. The 200-sample dataset I constructed specifically reflects the mix of Reddit posts, tweets, and news articles that the system actually processes, which I felt was more honest than reporting accuracy on a cleaner, professionally annotated corpus.

\section{Aspect-Based Sentiment Analysis}

Standard sentiment classification assigns a single polarity label to an entire document or sentence, discarding potentially contradictory signals within the same text.  An earnings call transcript, for instance, may express optimism about revenue growth while acknowledging concern about debt levels.  Aspect-Based Sentiment Analysis (ABSA) \citep{pontiki2014semeval} addresses this by identifying sentiment expressed towards specific \textit{aspects} or entities within text, enabling a finer-grained representation.

Pontiki et al.\ \citep{pontiki2014semeval} formalised ABSA through the SemEval shared tasks, establishing benchmark datasets covering aspect category detection, aspect term extraction, and aspect polarity classification.  While full ABSA systems require fine-tuned sequence labelling models and are computationally demanding, lighter-weight approximations using keyword-triggered sentence extraction have been proposed for high-throughput applications where approximate granularity suffices.

In financial text, aspect-level analysis matters because investor decisions often hinge on a specific facet of a company (its revenue trajectory, debt position, product pipeline, or macro exposure) rather than its sentiment overall. That observation drove me to build a lightweight ABSA module rather than simply reporting a headline positive/negative/neutral score. I wanted users to be able to see \textit{why} the model arrived at a given label, not just what that label was.

\section{Explainable AI in Financial Systems}

Traditional machine learning models in finance operate as black boxes, producing outputs without exposing the reasoning that generated them.  Doshi-Velez and Kim \citep{doshivelez2017interpretable} formalise interpretability as the ability to explain or present model behaviour in understandable terms to a human, arguing that interpretability requirements arise wherever model objectives cannot be fully specified in advance.  In financial applications, this is almost always the case: trust, auditability, and regulatory compliance all demand justification that a single confidence score cannot provide \citep{yeo2023xai}.

LIME \citep{ribeiro2016lime} explains individual predictions by fitting a locally interpretable linear surrogate model around a single input.  For text classification, LIME systematically masks or removes words, queries the black-box model on the perturbed inputs, and fits a sparse linear model whose weights identify which words most influenced the original prediction.  LIME is model-agnostic, computationally practical, and produces human-readable explanations, making it well-suited for embedding within user-facing interfaces.

SHAP \citep{lundberg2017shap} extends interpretability using Shapley values from cooperative game theory, guaranteeing consistency and local accuracy properties that LIME does not formally provide.  In NLP tasks, SHAP assigns per-token contributions to a model's output, making it attractive for corpus-level auditing and identifying systematically influential vocabulary.  The higher computational cost of SHAP relative to LIME makes it more appropriate for offline analysis and notebook exploration than for real-time API endpoints.

Yeo et al.\ \citep{yeo2023xai} survey the state of XAI in financial applications, finding that most current work focuses on model-level interpretability rather than integrating explanations within interactive decision-support tools.  Recent work by \citeauthor{bala2025integration} (\citeyear{bala2025integration}) demonstrates that XAI integration with deep learning can support real-time sentiment-driven stock price prediction, reinforcing the importance of transparent architectures in production-grade financial AI.  I designed the system specifically to close this gap by embedding LIME explanations directly within the dashboard's frontend, so that a user can see token-level reasoning without ever leaving the interface.

\section{Sentiment--Price Relationships and Causal Inference}

Beyond demonstrating correlation between sentiment and returns, a strand of the literature investigates the direction and timing of that relationship.  Simple Pearson and Spearman correlation coefficients capture the strength and monotonicity of co-movement, but they do not speak to causality or to the temporal structure of information propagation.

Lag analysis, computing correlation between sentiment at time \(t\) and price at time \(t+k\) for varying lags \(k\), is a standard preliminary step for identifying whether sentiment leads price movements, and, if so, by how many days.  Schumaker and Chen \citep{schumaker2009textmining} demonstrated that financial news sentiment predicts short-term price movements with lags of one to two hours in intraday data, while social media data has been associated with predictive horizons of one to three days at the daily granularity \citep{bollen2011twitter}.

Granger causality testing \citep{granger1969} provides a more formal framework for investigating temporal precedence.  A variable $X$ Granger-causes $Y$ if past values of $X$ contain statistically significant predictive information about future values of $Y$, above and beyond the past values of $Y$ alone.  Granger causality has been widely applied in macroeconomics and finance to study information transmission between markets.  For sentiment--price analysis, it offers a principled method to test whether aggregate daily sentiment scores carry incremental predictive information about next-day returns, without making the stronger claim of structural causation.

Rolling correlation further enriches the picture by tracking how the sentiment--price relationship evolves through time, revealing regimes in which the two signals co-move strongly and periods in which they decouple, for instance during periods of high market dislocation where price is driven by mechanical factors rather than sentiment.

The stationarity requirements of Granger testing mean it should ideally be paired with cointegration analysis \citep{engle1987cointegration} to rule out spurious regressions, a detail that many applied sentiment studies overlook.  Despite the maturity of these methods individually, I found no prior system that combined them all within a single interactive dashboard where the user can select a ticker, adjust the date range, and compare Pearson, Spearman, Granger, and rolling correlation side by side. Building that was one of the parts of this project I found most technically satisfying, partly because Granger causality is genuinely non-trivial to implement correctly.

\section{Gaps Identified}

Having reviewed this literature, I identified several recurring gaps that directly shaped what I chose to build:

\begin{itemize}
    \item \textbf{Single-source limitations.}
    The foundational studies in this area each rely on one platform:
    Tetlock \citeyearpar{tetlock2007giving} used \textit{Wall Street Journal}
    columns, Bollen et al.\ \citeyearpar{bollen2011twitter} used Twitter,
    Schumaker and Chen \citeyearpar{schumaker2009textmining} used financial
    news only, and Kraaijeveld and De Smedt \citeyearpar{kraaijeveld2020predictive}
    used Twitter for cryptocurrency.
    Long et al.\ \citeyearpar{long2023reddit} show that Reddit retail
    sentiment can diverge substantially from these channels during the same
    event, demonstrating that single-source signals are structurally
    incomplete. Multi-source aggregation with source-tagged outputs and
    comparable preprocessing remains underexplored.

    \item \textbf{Coarse sentiment representation.}
    The standard benchmark for financial sentiment, Financial PhraseBank
    \citep{malo2014financialphrasebank}, and the dominant model trained on
    it, FinBERT \citep{araci2019finbert}, both operate on three classes.
    Bollen et al.\ \citeyearpar{bollen2011twitter} proposed multi-dimensional
    mood indices for general Twitter data, but no finance-specific emotion
    taxonomy grounded in investor behavioural patterns (fear, scepticism,
    confidence, etc.) has been proposed in the literature.

    \item \textbf{Absence of lightweight ABSA for enrichment.}
    The SemEval ABSA shared tasks \citep{pontiki2014semeval} established
    full ABSA as a resource-intensive sequence-labelling problem requiring
    annotated training data.  Lightweight, keyword-driven evidence extraction
    for financial aspects in high-throughput ingestion pipelines has not been
    described; the gap between full ABSA and no aspect analysis at all is
    unaddressed.

    \item \textbf{XAI confined to offline tools.}
    \citeauthor{yeo2023xai} (\citeyear{yeo2023xai}) survey XAI in financial
    applications and explicitly find that deployment within interactive
    user-facing tools is largely absent.  \citeauthor{bala2025integration}
    (\citeyear{bala2025integration}) demonstrate that deep learning with XAI
    can support real-time stock prediction, but still do not integrate
    explanations into a production frontend that non-technical users can
    access directly.

    \item \textbf{Fragmented evaluation.}
    Published studies typically report top-line accuracy or F1 on
    Financial PhraseBank \citep{malo2014financialphrasebank}, without
    per-class breakdowns, confusion matrix analysis, qualitative XAI
    inspection, or system-level performance measurements.
    This makes it difficult to understand \textit{how} a model fails, not
    just whether it does.

    \item \textbf{No integrated correlation--causality dashboard.}
    Bollen et al.\ \citeyearpar{bollen2011twitter}, Schumaker and Chen
    \citeyearpar{schumaker2009textmining}, and Tetlock
    \citeyearpar{tetlock2007giving} each demonstrate correlation or causal
    relationships in standalone analyses.  No prior work describes an
    interactive, ticker-agnostic tool combining Pearson, Spearman, lag
    analysis, Granger causality, rolling correlation, and market adjustment
    within a single user-configurable interface.
\end{itemize}

\section{Summary}

Overall, the literature establishes that investor sentiment derived from multi-platform text influences market behaviour and can be captured at scale using domain-adaptive models like FinBERT. XAI techniques address the interpretability gap, and Granger causality provides a principled framework for testing sentiment--price dynamics. ABSA points toward more granular analysis of what investors are actually discussing. What I did not find was any single system that combines all of these elements (multi-source ingestion, enriched emotion taxonomy, lightweight ABSA, production XAI, and an interactive correlation engine) into one reproducible, user-facing tool. Addressing that gap is the core purpose of this project.
